\documentclass[11pt]{article}
\usepackage[margin=1in]{geometry}
\usepackage{amsmath,amssymb,amsthm}
\usepackage{bm}
\usepackage{booktabs}
\usepackage{hyperref}
\hypersetup{colorlinks=true,linkcolor=blue,urlcolor=blue}

\newcommand{\R}{\mathbf{R}}
\newcommand{\D}{\mathbf{D}}
\newcommand{\K}{\mathbf{K}}
\newcommand{\I}{\mathbf{I}}
\newcommand{\bet}{{\bm{\beta}}}
\newcommand{\al}{{\bm{\alpha}}}
\newcommand{\Sal}{{\bm{\Sigma}_{\bm{\alpha}}}}
\newcommand{\Ai}{A_i}
\newcommand{\E}{\mathbb{E}}
\newcommand{\Var}{\operatorname{Var}}
\newcommand{\Cov}{\operatorname{Cov}}
\newcommand{\N}{\mathcal{N}}

\title{\textbf{Spatial confounding and misaligned covariates in SDALGCP2}\\[4pt]
\large Restricted spatial regression, and covariates measured on a different
support from the outcome}
\author{SDALGCP2 development notes}
\date{\today}

\begin{document}
\maketitle

\begin{abstract}
Two practical problems in areal disease mapping with spatially structured
predictors. First, a spatially smooth covariate is collinear with the spatial
random effect, which absorbs its signal and biases the regression coefficient
(\emph{spatial confounding}); we describe the restricted spatial regression
implemented in \texttt{SDALGCP2}. Second, the covariate may be observed on a
\emph{different spatial support} from the outcome -- other polygons, point
monitors, or a coarser grid; we derive a change-of-support treatment that predicts
the covariate to the candidate points feeding the intensity-scale offset and
propagates the prediction uncertainty through a Berkson correction.
\end{abstract}

\section{Background: the SDA-LGCP model}

For aggregated counts $Y_i$ over regions $\Ai$, $i=1,\dots,N$, the SDA-LGCP sets
\begin{equation}
Y_i \mid \bm\eta \sim \operatorname{Poisson}(m_i e^{\eta_i}),
\qquad \bm\eta = \D\bet + \mathbf{Z},
\qquad \mathbf{Z}\sim\N(\mathbf 0,\ \sigma^2\R(\phi)),
\label{eq:model}
\end{equation}
where $\D$ ($N\times p$) is the region-level design, $m_i$ an offset, and
$\R(\phi)$ the aggregated spatial correlation built from candidate points inside
each region.

\section{Part I: spatial confounding}

\subsection{The problem}
When a covariate column of $\D$ is spatially smooth, it lies close to the span of
the leading modes of $\R(\phi)$. The random effect $\mathbf Z$ can then reproduce
much of that column, so the data are equally well explained by ``covariate effect''
or by ``spatial noise''. The fitted $\widehat\bet$ is pulled away from the value a
non-spatial GLM would give, usually \emph{attenuated} toward zero, with inflated
variance. In \texttt{SDALGCP2} on a simulated gradient covariate (true
$\beta=0.5$) the standard fit returns $\widehat\beta\approx0.28$ while a Poisson GLM
returns $\approx0.42$ -- the spatial term has eaten the signal.

\subsection{Restricted spatial regression (RSR)}
Following Reich, Hodges \& Zadnik (2006) and Hughes \& Haran (2013), constrain the
random effect to the orthogonal complement of the fixed-effect design. Let $\K$
($N\times r$, $r=N-p$) be an orthonormal basis of $\operatorname{null}(\D^\top)$, so
\begin{equation}
\K^\top\K=\I_r,\qquad \D^\top\K=\mathbf 0 .
\end{equation}
(\texttt{SDALGCP2} takes $\K$ from the trailing columns of a complete QR of $\D$.)
Replace the unrestricted random effect by $\K\al$:
\begin{equation}
\boxed{\;
Y_i\mid\al\sim\operatorname{Poisson}(m_i e^{\eta_i}),\quad
\bm\eta=\D\bet+\K\al,\quad
\al\sim\N\!\big(\mathbf 0,\ \sigma^2\,\K^\top\R(\phi)\K\big).\;}
\label{eq:rsr}
\end{equation}
Because $\K\al\in\operatorname{col}(\D)^\perp$, the random effect \emph{cannot}
reproduce any column of $\D$; $\bet$ is identified by the data, not by the spatial
structure, and the spatial term only smooths the residual. RSR removes the
confounding by construction; in the Gaussian case $\widehat\bet$ equals the OLS
estimate, and for the Poisson model it is close to the GLM estimate.

\subsection{Estimation: Laplace-approximate marginal likelihood}
Rather than Monte Carlo, \eqref{eq:rsr} is fitted by integrating $\al$ out at its
mode. With $\Sal(\phi)=\K^\top\R(\phi)\K$ and
$g(\al)=\sum_i\{y_i\eta_i-m_i e^{\eta_i}\}-\tfrac1{2\sigma^2}\al^\top\Sal^{-1}\al$,
let $\widehat\al=\arg\max_\al g(\al)$ (a few Newton steps; the negative Hessian is
$\mathbf H=\K^\top\operatorname{diag}(m_i e^{\eta_i})\K+\Sal^{-1}/\sigma^2$).
The Laplace-approximate marginal log-likelihood is
\begin{equation}
\ell(\bet,\sigma^2,\phi)=
\sum_i\{y_i\widehat\eta_i-m_i e^{\widehat\eta_i}\}
-\tfrac12\Big(r\log\sigma^2+\log|\Sal|
+\tfrac{\widehat\al^\top\Sal^{-1}\widehat\al}{\sigma^2}\Big)
-\tfrac12\log|\mathbf H| + \text{const},
\label{eq:laplace}
\end{equation}
maximised over $(\bet,\log\sigma^2)$ for each $\phi$ on a grid, with standard errors
from the Hessian of \eqref{eq:laplace}. On the simulation above RSR returns
$\widehat\beta\approx0.44$ (SE $0.07$) -- de-confounded -- while still estimating the
spatial parameters. Use \texttt{sdalgcp(..., control = sdalgcp\_control(confounding
= "restricted"))}.

\subsection{Spatio-temporal restricted regression}
With the same regions observed at times $t=1,\dots,T$, stack the $NT$ region--times
and let $\D$ ($NT\times p$) be the space--time fixed-effect design and $\K$
($NT\times(NT-p)$) an orthonormal basis of $\operatorname{null}(\D^\top)$. The
restricted space--time model keeps \eqref{eq:rsr} but with the \emph{separable}
space--time covariance,
\begin{equation}
\bm\eta=\D\bet+\K\al,\qquad
\al\sim\N\!\Big(\mathbf 0,\ \sigma^2\,\K^\top\big(\R_t(\nu)\otimes\R_s(\phi)\big)\K\Big),
\label{eq:rsr-st}
\end{equation}
where $\R_s(\phi)$ is the aggregated spatial correlation and $\R_t(\nu)$ a temporal
Mat\'ern correlation. Because $\K\al\in\operatorname{col}(\D)^\perp$ over all
region--times, no spatial \emph{or} temporal mode of the random effect can
reproduce a column of $\D$, so $\bet$ is de-confounded exactly as before.
Estimation is the same Laplace-approximate marginal likelihood \eqref{eq:laplace}
with $\Sal=\K^\top(\R_t(\nu)\otimes\R_s(\phi))\K$, profiled over the spatial grid
$\phi$ and maximised over the temporal range $\nu$ alongside $(\bet,\sigma^2)$. The
construction reduces \emph{exactly} to the spatial RSR \eqref{eq:rsr} at $T=1$ (then
$\R_t\equiv1$ and $\nu$ drops out), a check used in the package tests. It forms the
full $NT\times NT$ covariance, so it is $O\!\big((NT)^3\big)$ and best suited to
modest $NT$. Use \texttt{SDALGCP2\_ST(..., confounding = "restricted")}.

\section{Part II: covariates on a different support}

\subsection{The problem}
The intensity-scale treatment of a continuous covariate needs its value
$z(x_{ik})$ at the candidate points $x_{ik}$ inside each outcome region, entering
the offset
$b_i(\bet)=\log\sum_k w_{ik}\exp\{z(x_{ik})^\top\bet\}$
(companion note on raster covariates). A raster gives $z$ everywhere, so this is
immediate. But often the covariate is observed on a \emph{different support}:
\begin{enumerate}
\item \textbf{point} locations $s_1,\dots,s_L$ (air-quality monitors, weather
  stations);
\item \textbf{other polygons} $B_1,\dots,B_M$ (census tracts, catchments) that do
  not nest with the outcome regions $\Ai$;
\item a \textbf{coarser/finer grid} than the candidate points.
\end{enumerate}
All three are \emph{spatial misalignment} (change-of-support) problems
(Gotway \& Young, 2002; Gelfand, 2010): we must transfer $z$ to the points
$x_{ik}$.

\subsection{Predict the covariate, then aggregate}
Model the covariate as a latent spatial process $z(\cdot)$ and predict it to the
candidate points from the observed data $\mathbf z^{\mathrm{obs}}$.

\paragraph{Point support.} Fit a geostatistical (Gaussian-process) model to
$\{z(s_l)\}$ and krige to each $x_{ik}$, giving a predictive mean and variance
\begin{equation}
\mu_z(x_{ik})=\E[z(x_{ik})\mid \mathbf z^{\mathrm{obs}}],
\qquad
v_z(x_{ik})=\Var[z(x_{ik})\mid \mathbf z^{\mathrm{obs}}].
\end{equation}

\paragraph{Areal support.} If $z$ is observed as averages over polygons $B_j$,
$z^{\mathrm{obs}}_j=|B_j|^{-1}\!\int_{B_j} z(u)\,du$, model the underlying surface
$z(\cdot)$ and use \emph{areal kriging}: the observations are linear (integral)
functionals of $z$, so the joint Gaussian model gives the same predictive
$\mu_z,v_z$ at the points $x_{ik}$ -- exactly the discrete-aggregation machinery
\texttt{SDALGCP2} already uses, applied to the covariate.

\paragraph{Grid support.} A special case of point/areal kriging; bilinear
interpolation is the cheap plug-in.

\subsection{Berkson correction: propagate the prediction uncertainty}
Plugging the kriged mean $\mu_z(x_{ik})$ straight into $b_i(\bet)$ ignores that
$z(x_{ik})$ is \emph{predicted}, not observed -- a classic errors-in-variables
problem that attenuates $\bet$. Because the kriging predictive distribution is
Gaussian, $z(x_{ik})\mid\mathbf z^{\mathrm{obs}}\sim\N(\mu_z,v_z)$, we can integrate
it out in closed form. Using $\E[e^{a^\top X}]=\exp(a^\top\mu+\tfrac12 a^\top V a)$
for $X\sim\N(\mu,V)$,
\begin{equation}
\E\big[\exp\{z(x_{ik})^\top\bet\}\,\big|\,\mathbf z^{\mathrm{obs}}\big]
=\exp\!\Big\{\mu_z(x_{ik})^\top\bet+\tfrac12\,\bet^\top V_z(x_{ik})\,\bet\Big\},
\end{equation}
so the offset that correctly accounts for the misalignment uncertainty is
\begin{equation}
\boxed{\;
b_i(\bet)=\log\sum_k w_{ik}\,
\exp\!\Big\{\mu_z(x_{ik})^\top\bet+\tfrac12\,\bet^\top V_z(x_{ik})\,\bet\Big\}.\;}
\label{eq:berkson}
\end{equation}
Equation \eqref{eq:berkson} is the ordinary intensity-scale offset evaluated at the
predicted covariate mean, \emph{plus} a $\tfrac12\bet^\top V_z\bet$ term that inflates
the contribution of poorly-observed locations -- the Berkson correction. It reduces
to the raster case when $V_z\to 0$ (the covariate is known), and it is fitted by
the same Gauss--Newton scheme as the raster model, with the effective covariate
gradient $\partial b_i/\partial\bet$ now including the $V_z\bet$ term. Full
propagation of the covariate model's own parameter uncertainty (a two-stage or
joint model) is a natural further step.

\paragraph{Summary.} Misaligned covariates are handled in two steps that reuse the
existing machinery: (i) predict $z$ to the candidate points by (areal) kriging,
obtaining $\mu_z,v_z$; (ii) feed them into the Berkson-corrected offset
\eqref{eq:berkson}. The candidate-point grid that already discretises the outcome
regions doubles as the common support on which outcome and covariate are aligned.

\paragraph{Spatio-temporal misalignment.} When the outcome is observed over several
times but the covariate support is time-invariant (e.g.\ fixed monitors), the
covariate is kriged once to the candidate points and the Berkson offset
\eqref{eq:berkson} feeds the Kronecker-free space--time fit exactly as in the
spatial case, inside the same Gauss--Newton loop
(\texttt{SDALGCP2\_ST(..., covariates = )}).

\paragraph{References.}
Reich, B.J., Hodges, J.S., Zadnik, V. (2006). Effects of residual smoothing on the
posterior of the fixed effects in disease-mapping models. \emph{Biometrics} 62,
1197--1206.\\
Hughes, J., Haran, M. (2013). Dimension reduction and alleviation of confounding
for spatial generalized linear mixed models. \emph{JRSS-B} 75, 139--159.\\
Gotway, C.A., Young, L.J. (2002). Combining incompatible spatial data.
\emph{JASA} 97, 632--648.
\end{document}
